\begin{abstract}
We study a weighted inert-versus-split prime race in the one-parameter hyperelliptic
pencil
\(y^2=\prod_{a=1}^{2g}(x-a)(x-t)\) over finite fields.  For each curve, the
normalized race sequence has a Ces\`aro endpoint law.  In the critical
scaling \(m+1/2=\lambda/((2g+1)\log Q)\), we prove that these quenched laws,
viewed as random elements of \(\mathcal P_2(\mathbb R)\), converge to the
random law generated by a marked symplectic hard-edge process.  The limit is
simultaneous in genus and field size for odd square prime powers satisfying
\(\log Q/(g^2\log(g+2))\to\infty\), and we give a completely specified
sequence of fields.  The limiting random measure and its positive-half-line
mass are both nonconstant.  The proof combines density-one linear
independence, genus-explicit Frobenius equidistribution, heat-kernel smoothing,
and direct \(W_2\) control of the marked series.  This is a
large-field/large-genus bridge theorem, not a fixed-field microscopic result.
\end{abstract}

\section{Introduction and statement of the theorem}
\label{sec:introduction}

Prime races over a function field turn the zeros of an arithmetic
\(L\)-function into a finite collection of Frobenius angles.  For a fixed
curve, linear independence of those angles converts the endpoint variable
into a random-phase cosine sum; for a family of curves, monodromy and
equidistribution make the coefficients themselves random.  These two kinds of
randomness interact most sharply at the hard edge, where the weight is tuned
so that the lowest Frobenius angles contribute at order one.

This paper isolates that interaction in one explicit family and one
simultaneous limit.  Fix \(\lambda>0\).  For each \(g\ge2\), let \(Q_g\) be an
odd square prime power of characteristic \(p_g>2g+1\), put
\[
 f_g(x)=\prod_{a=1}^{2g}(x-a),\qquad
 U_g=\mathbb A^1\setminus\{1,\ldots,2g\},
\]
and consider the pencil
\begin{equation}
 C_{g,t}:y^2=f_g(x)(x-t),\qquad t\in U_g(\F_{Q_g}).
\label{eq:pencil}
\end{equation}
The curve has genus \(g\).  Its quadratic character \(\chi_{g,t}\) is
\(+1\) at split prime polynomials and \(-1\) at inert ones.  Write
\begin{equation}
 d_g=2g+1,\qquad
 M_g=m_g+\frac12=\frac{\lambda}{d_g\log Q_g}.
\label{eq:critical-scale}
\end{equation}
For an endpoint \(N\), the normalized inert-minus-split race is
\begin{equation}
 \mathcal E_{g,t,\lambda}(N)
 =2(1-Q_g^{-M_g})Q_g^{-M_gN}
 \left(-\sum_{\substack{P\ \mathrm{monic\ irreducible}\\ \deg P\le N}}
       (\deg P)\chi_{g,t}(P)Q_g^{m_g\deg P}\right).
\label{eq:race}
\end{equation}

Write \(\Ptwo\) for the space of Borel probability measures on \(\mathbb R\)
with finite second moment, equipped with the quadratic Wasserstein metric
\(\Wtwo\).

For every parameter \(t\), the empirical endpoint distribution of
\(\mathcal E_{g,t,\lambda}(N)\) exists; this follows from
Kronecker--Weyl on the closure of the cyclic angle--parity orbit.  We denote
the genuine endpoint law by \(\nu_{g,t,\lambda}\in\Ptwo\).  When the positive
Frobenius angles together with \(\pi\) are linearly independent over
\(\mathbb Q\), that orbit is the full phase torus with an independent fair
parity coordinate, giving the explicit law in \eqref{eq:finite-quenched-law}.

The limiting object is a random probability measure.  Let \(\PiSp\) be the
determinantal process on \([0,\infty)\) with kernel
\begin{equation}
 K_{\mathrm{Sp}}(x,y)
 =\frac{\sin\pi(x-y)}{\pi(x-y)}
  -\frac{\sin\pi(x+y)}{\pi(x+y)}.
\label{eq:sp-kernel}
\end{equation}
For a locally finite configuration \(\xi\), define
\begin{equation}
 a_\lambda(y)=\frac{4\lambda}
 {\sqrt{\lambda^2+(2\pi y)^2}},\qquad
 \nu_\lambda(\xi)
 =\Law_\Phi\left(1+\sum_{y\in\xi}a_\lambda(y)\cos\Phi_y\right),
\label{eq:quenched-limit-law}
\end{equation}
where, conditionally on \(\xi\), the phases are independent and uniform on
\([0,2\pi]\).  The series converges in conditional \(L^2\) almost surely.
There are now two probability layers.  The phases are integrated out to form
the conditional, or quenched, measure; the uniform curve parameter \(t\), and
in the limit the point configuration \(\PiSp\), supply the outer randomness
that makes this measure itself random.  Thus \(\Longrightarrow\) below means
weak convergence of the outer laws on the Polish space \((\Ptwo,\Wtwo)\).
The coefficient sequences are naturally square-summable, and the shared-phase
coupling controls their second moments exactly, which is why \(\Wtwo\) is the
topology carried by the proof.

\begin{theorem}[Simultaneous quenched critical law]
\label{thm:main}
For every fixed \(\lambda>0\), suppose
\begin{equation}
 \frac{\log Q_g}{g^2\log(g+2)}\longrightarrow\infty.
\label{eq:growth}
\end{equation}
Choose \(t_g\) uniformly from \(U_g(\F_{Q_g})\).  Then, as random elements of
\((\Ptwo,\Wtwo)\),
\begin{equation}
 \nu_{g,t_g,\lambda}\Longrightarrow\nu_\lambda(\PiSp).
\label{eq:main-convergence}
\end{equation}
The limiting random probability measure is not almost surely constant and is
almost surely absolutely continuous.
\end{theorem}

\begin{corollary}[Prime-race densities]
\label{cor:prime-race-densities}
Under the hypotheses of \cref{thm:main}, the positive-half-line masses satisfy
\begin{equation}
 \nu_{g,t_g,\lambda}((0,\infty))
 \Longrightarrow
 \nu_\lambda(\PiSp)((0,\infty)),
\label{eq:density-convergence}
\end{equation}
and
\[
 \lim_{g\to\infty}\E_{t_g}\,
 \nu_{g,t_g,\lambda}((0,\infty))
 =\E\,\nu_\lambda(\PiSp)((0,\infty)).
\]
The limiting scalar \(\nu_\lambda(\PiSp)((0,\infty))\) is not almost surely
constant.  With probability \(1-o(1)\), the left side of
\eqref{eq:density-convergence} is the natural density of positive endpoints in
the original race.
\end{corollary}

\begin{example}[An explicit admissible sequence]
\label{ex:explicit-fields}
The conclusions of \cref{thm:main,cor:prime-race-densities} hold when \(p_g\)
is the least prime greater
than \(2g+1\) and
\begin{equation}
 Q_g=p_g^{2n_g},\qquad
 n_g=\min\left\{n\ge1:p_g^{2n}\ge
 \exp\!\bigl(g^2\log^2(g+2)\bigr)\right\}.
\label{eq:explicit-fields}
\end{equation}
Indeed, \cref{eq:explicit-field-growth} verifies \cref{eq:growth}.
\end{example}

The order of limits in \cref{thm:main} is not iterated: the genus, the
characteristic, and the extension degree vary along one specified sequence.
The square-field hypothesis is structural.  It makes \(\sqrt{Q_g}\) rational,
so that a maximal Weyl-group Galois condition on the Frobenius polynomial
implies the angle-plus-\(\pi\) independence required by the endpoint race.

The result is nevertheless not a fixed-field microscopic theorem.  Condition
\eqref{eq:growth} forces the field to grow very rapidly and lets quantitative
Frobenius equidistribution dominate all genus-dependent smoothing costs.  The
fixed-\(Q\), large-genus hard-edge problem remains outside the argument.

\paragraph{Relation to earlier work.}
Cha developed limiting distributions and Chebyshev bias for function-field
prime races~\cite{Cha2008}, with subsequent refinements and related family
results in~\cite{ChaKim2010,ChaFiorilliJouve2016,Bailleul2022}.  A
particularly close direct race precedent is Bailleul--Devin--Keliher--Li
\cite{BailleulDevinKeliherLi2024}: they study the same hyperelliptic
split-versus-inert setting, fixed-curve endpoint laws, parity centers, and the
rarity of exceptional biases, but not a moving critical weight or a random
measure-valued hard-edge limit.  Recent work on ties and on high-trace
hyperelliptic statistics further shows that neither function-field race
phenomena nor simultaneous large-field/large-genus information is new by
itself~\cite{BatesEtAl2026,LalinEtAl2026}.

Earlier race work therefore supplies fixed-curve scalar laws and bias
statistics, while hyperelliptic-family work supplies Frobenius
equidistribution and low-zero statistics.  The object studied here retains the
Frobenius configuration as outer randomness and proves convergence of the
conditional law itself in \(\Wtwo\), at the moving critical weight for which
the lowest angles contribute at order one.  The three bridges established are the
exact weighted-race reduction, a genus-explicit transfer for a nonsmooth
measure-valued test, and closed-edge \(\ell^2\) stability of the marked law.

The
combination of generic Frobenius linear independence with family averaging
also appears in Cha's hyperelliptic M\"obius study and in Perret--Gentil's
trace-function families~\cite{ChaMobius2010,PerretGentil2019}; their
statistics and limiting objects differ from the quenched critical law here.
The
hyperelliptic monodromy and quantitative equidistribution inputs used here are
due to Katz--Sarnak~\cite{KatzSarnak1999}; density-one maximal Galois group and
the relation-space criterion are due to Kowalski
~\cite{Kowalski2006,KowalskiErratum,Kowalski2008}.  Low-zero statistics for
hyperelliptic families have been studied in several regimes, including
~\cite{EntinRodittyGershonRudnick2013,EntinPirani2024}.  Our contribution is
the critical weighted-race calculation and the quantitative passage from the
actual endpoint law to a nondegenerate random \(\Ptwo\)-valued hard-edge law
in one simultaneous arithmetic regime.  We do not claim that simultaneous
large field and large genus, by themselves, are new.

\paragraph{Proof structure.}
The argument follows one chain:
\[
 \begin{aligned}
 \text{endpoint laws}
   &\xrightarrow{\mathrm{LI}}\text{Frobenius-phase laws}\\
   &\xrightarrow{\mathrm{heat\ bridge}}\text{Haar }\USp(2g)\text{ laws}\\
   &\xrightarrow{\mathrm{marked}\ \Wtwo\ \mathrm{stability}}
     \text{hard-edge quenched law}.
 \end{aligned}
\]
\Cref{sec:race-law} derives the exact random-phase law, including the two
prime-square parity centers.  \Cref{sec:linear-independence} proves that the
required linear independence fails on a set of density \(o(1)\).
\Cref{sec:transfer} transfers the probability-measure-valued Frobenius
functional to Haar \(\USp(2g)\), with every genus dependence visible.
\Cref{sec:hard-edge} proves the closed-hard-edge marked-law limit and both
nondegeneracy assertions.  \Cref{sec:assembly} assembles the estimates.  The
two analytic tools most sensitive to normalization are proved separately in
the appendices.
